% Appendix B, part two (Chapters 11-19). Included by app_b_solutions.tex via \input.

\section{Maximum Likelihood}
\noindent{\small Chapter~\ref{ch:mle}.}

\begin{solution}{11-1}
$\hat\lambda=\bar X$, and by invariance $\widehat{e^{-\lambda}}=e^{-\bar X}$. Since $T=n\bar X\sim\text{Poisson}(n\lambda)$, $\E e^{-\bar X}=\E e^{-T/n}=\exp\{n\lambda(e^{-1/n}-1)\}$. Since $e^{-1/n}-1>-1/n$, this value exceeds $e^{-\lambda}$: the MLE is biased upwards (a consequence of Jensen). For $\lambda=2$, $n=10$: $0.1491$ vs $0.1353$. The UMVUE $(1-1/n)^T$ is unbiased; in the simulation of Chapter~\ref{ch:yeterlilik} its MSE was also smaller.
\end{solution}

\begin{solution}{11-2}
$L=\lambda^ne^{-\lambda\sum(x_i-\mu)}\I\{\mu\le x_{(1)}\}$ is increasing in $\mu$ for fixed $\lambda$; $\hat\mu=X_{(1)}$. Then $\partial_\lambda\ell=n/\lambda-\sum(x_i-\hat\mu)=0$ gives $\hat\lambda=1/(\bar X-X_{(1)})$. $X_{(1)}-\mu\sim\text{Exp}(n\lambda)$ (Exercise~\ref{exr:04-3}), so $\Bias(\hat\mu)=\frac{1}{n\lambda}$. Since $\E[\bar X-X_{(1)}]=\frac1\lambda-\frac{1}{n\lambda}=\frac{n-1}{n\lambda}$,
\[
\tilde\mu=X_{(1)}-\frac{\bar X-X_{(1)}}{n-1}
\]
is unbiased.
\end{solution}

\begin{solution}{11-3}
$\ell(\theta)=(2n_1+n_2)\ln\theta+(n_2+2n_3)\ln(1-\theta)+\text{const}$. $\ell'=0$ gives $\hat\theta=\frac{2n_1+n_2}{2n}$, i.e.\ the proportion of the first allele among the $2n$ alleles in the sample. $-\ell''=\frac{2n_1+n_2}{\theta^2}+\frac{n_2+2n_3}{(1-\theta)^2}$; with $\E(2n_1+n_2)=2n\theta$ we get $I_n(\theta)=\frac{2n}{\theta(1-\theta)}$: the information of $2n$ independent Bernoulli trials.
\end{solution}

\begin{solution}{11-4}
MLE: $\hat\mu_i=\bar X_i=\frac{X_{i1}+X_{i2}}2$ and $\hat\sigma^2=\frac{1}{2n}\sum_i\sum_j(X_{ij}-\bar X_i)^2=\frac{1}{4n}\sum_i(X_{i1}-X_{i2})^2$. Since $X_{i1}-X_{i2}\sim\Normal(0,2\sigma^2)$, $\E\hat\sigma^2=\sigma^2/2$ and by the SLLN $\hat\sigma^2\asto\sigma^2/2$: information does not accumulate, because each new pair brings a new nuisance parameter. A consistent estimator: $\tilde\sigma^2=\frac{1}{2n}\sum_i(X_{i1}-X_{i2})^2$. Simulation ($n=1000$, $\sigma^2=4$): $\hat\sigma^2=2.16$, $\tilde\sigma^2=4.31$.
\end{solution}

\section{MLE Asymptotics and the Delta Method}
\noindent{\small Chapter~\ref{ch:mle-asimptotik}.}

\begin{solution}{12-1}
$I_1(\lambda)=1/\lambda^2$; by Theorem~\ref{thm:mle-an}, $\sqrt n(\hat\lambda-\lambda)\dto\Normal(0,\lambda^2)$. Delta method: by the CLT $\sqrt n(\bar X-1/\lambda)\dto\Normal(0,1/\lambda^2)$, $g(\mu)=1/\mu$, $g'(1/\lambda)=-\lambda^2$; the variance is $\lambda^4/\lambda^2=\lambda^2$. For $\ln\hat\lambda$, $g'(\lambda)=1/\lambda$ and the asymptotic variance is $\lambda^2/\lambda^2=1$: the $\ln$ transformation is variance-stabilising for the exponential model.
\end{solution}

\begin{solution}{12-2}
$g(a,b)=b/a$, $\nabla g=(-\mu_Y/\mu_X^2,\ 1/\mu_X)$. With $R=\mu_Y/\mu_X$, Theorem~\ref{thm:delta-cok} gives
\[
\sqrt n\Big(\frac{\bar Y}{\bar X}-R\Big)\dto\Normal\Big(0,\ \frac{\sigma_Y^2-2R\sigma_{XY}+R^2\sigma_X^2}{\mu_X^2}\Big)=\Normal\Big(0,\frac{\Var(Y-RX)}{\mu_X^2}\Big).
\]
Positive correlation stabilises the ratio estimate; the ``ratio estimator'' of survey sampling exploits this.
\end{solution}

\begin{solution}{12-3}
Simulation with $4\times10^4$ replications:
\begin{center}\small
\begin{tabular}{@{}lccc@{}}\toprule
 & $n\Var(\text{median})$ & $n\Var(\text{one-step})$ & theory\\\midrule
$n=51$ & $2.62$ & $2.13$ & $\pi^2/4=2.47$; \ $1/I_1=2$\\
$n=201$ & $2.50$ & $2.03$ & \\\bottomrule
\end{tabular}
\end{center}
The one-step estimator reduces the asymptotic variance of the median, $\approx2.47$, to the Fisher bound $2$. General theory (Le Cam): a single Newton--Raphson step (with expected or observed information) from any $\sqrt n$-consistent starting point is asymptotically equivalent to the MLE; for multimodal likelihoods (Example~\ref{ex:cauchy-mle}) it is a safe way to avoid searching for the global maximum.
\end{solution}

\begin{solution}{12-4}
$z=\tanh^{-1}r$ is approximately $\Normal(\tanh^{-1}\rho,\,1/(n-3))$; the interval is $\tanh\big(z\pm1.96/\sqrt{n-3}\big)$. (Using $\sqrt{n-3}$ instead of $\sqrt n$ is a small-sample correction.) For $\rho=0.8$, $n=20$, $2\times10^4$ replications: the Fisher $z$ coverage is $0.951$, that of the direct interval $r\pm1.96(1-r^2)/\sqrt n$ is $0.893$. The direct interval ignores the skewed distribution of $r$ near the boundary of $[-1,1]$.
\end{solution}

